\documentclass[11pt,a4paper]{article}
\usepackage[margin=1.8cm]{geometry}
\usepackage{amsmath,amssymb,tabularx,array,needspace,fontspec}
\setmainfont{FreeSerif}
\setlength{\parindent}{0pt}
\setlength{\parskip}{0pt}
\renewcommand{\arraystretch}{1.15}
\pagestyle{plain}
% \sq{roman}{question}
\newcommand{\sq}[2]{\Needspace{3\baselineskip}\par\smallskip\noindent\hangindent=5.8em\hangafter=1\makebox[5.8em][l]{\textbf{#1.}}#2\par}

\begin{document}

\begin{center}
  {\Huge\bfseries The Hope Science Academy}\\[3pt]
  {\small Excellence is our standard not goal}\\[2pt]
  {\small 247-E-1, (Opposite LDA School) Johar Town, Lahore. PH: 0300-8194789, 0322-4157001}
\end{center}

\medskip
\noindent
\begin{tabularx}{\linewidth}{|X|X|X|X|}
\hline
\textbf{Student Name:} & \textbf{Roll Num:} & \textbf{Class:} 12th & \textbf{Paper Code:} \\ \hline
\textbf{Subject Name:} Principles of Accounting & \textbf{Time Allowed:} & \textbf{Total Marks:} 237 & \textbf{Exam Date:} 02-Nov-2026 \\ \hline
\multicolumn{4}{|l|}{\textbf{Exam Syllabus:} FULL BOOK} \\ \hline
\end{tabularx}

\bigskip
\noindent\textbf{Q1. Write short answers of the following questions.}\hfill\textbf{1$\times$237=237}

\sq{I}{What do you understand by single entry system?}
\sq{II}{What is double entry system?}
\sq{III}{Under single entry system what kind of book-keeping records are kept?}
\sq{IV}{Clearly distinguish between single entry and double entry system of accounts.}
\sq{V}{How many methods are available for the determination of profit or loss under single entry system?}
\sq{VI}{What do you understand by increased-net worth method of single entry system?}
\sq{VII}{What is statement of affairs?}
\sq{VIII}{Write the equations to calculate the value of adjusted capital (at the end) and profit or loss under single entry system.}
\sq{IX}{What are the disadvantages of single entry system? (Any three)}
\sq{X}{Distinguish between statement of affairs and balance sheet.}
\sq{XI}{What is meant by non-profit making concerns?}
\sq{XII}{What is “Receipts and payment account”?}
\sq{XIII}{What is “income and expenditure account”?}
\sq{XIV}{Write two differences between cash book and receipts and payments account.}
\sq{XV}{Write the two differences between “Receipts and payments account” and “income and expenditure account”.}
\sq{XVI}{What are the advantages to prepare “Receipts and payments account”?}
\sq{XVII}{How “Receipts and payments account” is prepared?}
\sq{XVIII}{How “income and expenditure account” is prepared?}
\sq{XIX}{What is the difference between donation and subscription?}
\sq{XX}{What is special subscription?}
\sq{XXI}{What is the admission fee or entrance fee?}
\sq{XXII}{What do you mean by Legacy?}
\sq{XXIII}{What is meant by life membership fee?}
\sq{XXIV}{What is capital fund?}
\sq{XXV}{How you will deal each of these while preparing income and expenditure account and Balance sheet? Legacies, Donation for specific purpose, Donations (purpose is not given), Life membership fees, Entrance fees, Endowment fund receipts, Receipts for tournament fund.}
\sq{XXVI}{What is consignment?}
\sq{XXVII}{How many parties are involved in consignment?}
\sq{XXVIII}{What do you mean by consignor and consignee in consignment?}
\sq{XXIX}{How many classification of consignment are?}
\sq{XXX}{What are consignment inward and consignment outward?}
\sq{XXXI}{What is the relationship between consignor and consignee?}
\sq{XXXII}{Write the three differences between consignment and sale.}
\sq{XXXIII}{Define the following with reference to consignment: Commission, Del commission, Overriding commission}
\sq{XXXIV}{What is an account sale?}
\sq{XXXV}{What do you understand by preforms invoice?}
\sq{XXXVI}{What is abnormal loss of stock?}
\sq{XXXVII}{What is normal loss of stock?}
\sq{XXXVIII}{What is a consignment account?}
\sq{XXXIX}{What is the accounting treatment for abnormal loss of stock and normal loss of stock?}
\sq{XL}{How the value of unsold stock with the consignee can be calculated?}
\sq{XLI}{What will be the entries in consignor and consignee books for goods sent on consignment?}
\sq{XLII}{Pass the entries for expenses paid by consignor in consignor and consignee books.}
\sq{XLIII}{What entries will have to be made for receiving advance from consignee in consignor and consignee books?}
\sq{XLIV}{Pass the entries if the consignee has accepted a bill of exchange as an advance in consignor and consignee books.}
\sq{XLV}{Show the entries for expenses paid by consignee in the consignor and consignee books.}
\sq{XLVI}{For goods sold by consignee, pass the entries in consignor and consignee books.}
\sq{XLVII}{Show the entries for goods sold on credit by consignee in consignor and consignee books.}
\sq{XLVIII}{What will be the entries for commission in consignor and consignee books?}
\sq{XLIX}{Pass the entries for del credited commission in consignor and consignee books.}
\sq{L}{Pass the entries for goods taken over by consignee for personal use in consignor and consignee books.}
\sq{LI}{Show the entries for payment of expense by consignee personally in consignor and consignee books.}
\sq{LII}{Pass the entries for bad debts if the consignee receives the del credited commission in consignor and consignee books.}
\sq{LIII}{Pass the entries for bad debts if the consignee does not receive the del credited commission in consignor and consignee books.}
\sq{LIV}{Show the entries for unsold stock with the consignee in consignor and consignee books.}
\sq{LV}{What will be the entry for abnormal loss of goods in consignor book, if it is not insured?}
\sq{LVI}{Pass the entries for final payment by consignee to consignor in cash or bank draft in consignor and consignee books.}
\sq{LVII}{What entries will have to be made for final payment by accepting a bill of exchange by consignee in consignor and consignee books?}
\sq{LVIII}{Pass the entries for profit earned on consignment in consignor and consignee books.}
\sq{LIX}{Show the entries for loss on consignment in consignor and consignee books.}
\sq{LX}{What will be the entry when the goods sent on consignment transferred to trading account in consignor book?}
\sq{LXI}{Show the formula to calculate the value of adjusted per unit cost in case of normal loss.}
\sq{LXII}{Show one specimen for account sales.}
\sq{LXIII}{Calculate the value of unsold stock in consignment from the following particulars. Unsold cases Rs.200 \& Rs.75 each; expenses paid by consignor Rs.550; unloading expenses paid by consignee Rs.200; total cases sent on consignment Rs.600}
\sq{LXIV}{Calculate the value of adjusted per unit cost in case of normal loss from the following particulars :\\[2pt]\hspace*{2em}\begin{tabular}{@{}ll@{}}
\quad (a) Original cost of goods consigned & Rs. 300 \\
\quad (b) Direct expenses & Rs. 5,0 \textit{[cut off in the source]} \\
\quad (c) Total units of goods & Rs. 1,0 \textit{[cut off in the source]} \\
\quad (d) Units sold & 20 \\
\end{tabular}}
\sq{LXV}{What is meant by a joint stock company?}
\sq{LXVI}{Write the five characteristics of Joint Stock Company.}
\sq{LXVII}{What is meant by that a company has a separate legal entity?}
\sq{LXVIII}{Why Joint Stock Company is called voluntary association?}
\sq{LXIX}{What is meant by perpetual existence of a joint stock company?}
\sq{LXX}{What do you know about the common seal of a Joint Stock Company?}
\sq{LXXI}{What is meant by limited liability in case of a Joint Stock Company?}
\sq{LXXII}{What is meant by transferability of shares in public company?}
\sq{LXXIII}{Is the ownership of company separate from management?}
\sq{LXXIV}{How many kinds of companies are under the company ordinance 1984?}
\sq{LXXV}{What is meant by company limited by shares?}
\sq{LXXVI}{How many kinds of company limited by shares?}
\sq{LXXVII}{Define private limited company and public limited company.}
\sq{LXXVIII}{What do you understand by company limited by guarantee?}
\sq{LXXIX}{What is unlimited company?}
\sq{LXXX}{Define association not for profit.}
\sq{LXXXI}{Write the steps for the formation of a public limited company.}
\sq{LXXXII}{Define the following: Directors, Promoters, Underwriters, shareholders}
\sq{LXXXIII}{What are the three-legal basic document of a public limited company?}
\sq{LXXXIV}{What is memorandum of association?}
\sq{LXXXV}{What is the purpose of memorandum of association?}
\sq{LXXXVI}{What is meant by articles of association in case of public limited company?}
\sq{LXXXVII}{Define prospectus.}
\sq{LXXXVIII}{What is the main object of prospectus?}
\sq{LXXXIX}{What is the difference between private limited company and public limited company?}
\sq{XC}{Distinguish between memorandum of association and articles of association.}
\sq{XCI}{What is underwriting commission?}
\sq{XCII}{Write two differences between partnership and Joint Stock Company.}
\sq{XCIII}{Define a share.}
\sq{XCIV}{What is share capital?}
\sq{XCV}{Define the following terms in relation to share capital: Authorized capital, Issued capital, Called up capital, Paid up capital, Reserve capital, Subscribed capital}
\sq{XCVI}{Distinguish between par value, book value and market value of shares.}
\sq{XCVII}{Define primary expenses/ preliminary expenses. What is the treatment of primary expenses in final accounts?}
\sq{XCVIII}{What do you understand by under subscription and over subscription in case of a joint stock company?}
\sq{XCIX}{What will be the entry for shares application money received?}
\sq{C}{What will be the entry for share application money refunded?}
\sq{CI}{What is meant by issue of shares at par?}
\sq{CII}{Pass the journal entry for allotment of shares at par.}
\sq{CIII}{What do you mean by the issue of shares at premium?}
\sq{CIV}{Pass the journal entry for allotment of shares at premium.}
\sq{CV}{What is meant by issue of shares at discount?}
\sq{CVI}{Pass the journal entry for allotment of shares at discount.}
\sq{CVII}{In case of under subscription, pass the entry for deficiency which is made up by the underwriter.}
\sq{CVIII}{Pass the journal entry for shares taken up by the promoters against preliminary expenses.}
\sq{CIX}{Pass the journal entry for shares taken up by the directors.}
\sq{CX}{What is the maximum discount rate issue of shares at discount?}
\sq{CXI}{What are the reasons for the issue of shares other than first or initial issue?}
\sq{CXII}{What is meant by issue of shares in consideration of assets other than cash?}
\sq{CXIII}{How premium on shares can be utilized under the legal provisions of company ordinance? (any two)}
\sq{CXIV}{Under what conditions a company can issue its shares at discount.}
\sq{CXV}{Pass the journal entry for allotment of shares to vendor against purchase consideration.}
\sq{CXVI}{Define debenture.}
\sq{CXVII}{Write the three names of the various kinds of debentures issued by a company.}
\sq{CXVIII}{Define the following different kinds of debentures: (a) Redeemable (b) Irredeemable (c) Simple or naked debentures (d) Registered debentures (e) Bearer debentures (f) Mortgage debentures (g) Debenture stock}
\sq{CXIX}{Distinguish between a share holder and a debenture holder. (any two)}
\sq{CXX}{Pass the journal entries for issue of debentures at par and redeemable at par.}
\sq{CXXI}{What journal entries will have to be made for issue of debentures at discount and redeemable at par?}
\sq{CXXII}{What journal entries will have to be made for issue of debentures at premium and redeemable at par?}
\sq{CXXIII}{Pass the journal entries for issue of debentures at par and redeemable at premium.}
\sq{CXXIV}{Pass the journal entries for issue of debentures at discount and redeemable at premium.}
\sq{CXXV}{Define the term depreciation.}
\sq{CXXVI}{How many types of depreciation are?}
\sq{CXXVII}{What are the causes of internal depreciation?}
\sq{CXXVIII}{What is meant by wear and tear in depreciation?}
\sq{CXXIX}{What do you mean by the term depletion?}
\sq{CXXX}{What are the causes of external depreciation?}
\sq{CXXXI}{What is obsolescence?}
\sq{CXXXII}{What is meant by efflux of time in depreciation?}
\sq{CXXXIII}{Why the provision of depreciation is necessary?}
\sq{CXXXIV}{What is the difference between depreciation and fluctuation in market price?}
\sq{CXXXV}{Write the three characteristics of depreciation.}
\sq{CXXXVI}{Write the three points for the ascertainment of the amount or rate of depreciation.}
\sq{CXXXVII}{What is meant by cost price of an asset?}
\sq{CXXXVIII}{What is meant by market price of an asset?}
\sq{CXXXIX}{What is scrap value of asset?}
\sq{CXL}{Define the term amortization.}
\sq{CXLI}{Distinguish between amortization.}
\sq{CXLII}{Distinguish between depreciation and amortization.}
\sq{CXLIII}{What is the treatment of depreciation in profit and loss account and in balance sheet?}
\sq{CXLIV}{What is the journal entry for depreciation?}
\sq{CXLV}{What do you mean by the term accumulated depreciation?}
\sq{CXLVI}{Define partnership.}
\sq{CXLVII}{What are the essential elements of partnership?}
\sq{CXLVIII}{What is the maximum number strength of partners in partnership?}
\sq{CXLIX}{What is partnership agreement?}
\sq{CL}{Write the four points of partnership agreement.}
\sq{CLI}{Write the three rights of partners in partnership.}
\sq{CLII}{Write three duties of partners in partnership.}
\sq{CLIII}{What are the rules, which are applicable in the absence of agreement in partnership? (any three)}
\sq{CLIV}{What is meant by fixed and fluctuating capitals in partnership?}
\sq{CLV}{Pass journal entry for capital introduced by different partners in partnership.}
\sq{CLVI}{What will be the journal entry for interest on capital in partnership?}
\sq{CLVII}{What will be the closing entry for interest on capital in partnership?}
\sq{CLVIII}{What will be the journal entry for salary payable to partners?}
\sq{CLIX}{What will be the closing entry, when the salary payable is transferred to profit and loss appropriation account?}
\sq{CLX}{What will be journal entry for drawings in cash in partnership?}
\sq{CLXI}{What will be the journal entry for interest on drawings in partnership?}
\sq{CLXII}{What will be the closing entry, when interest on drawings is transferred to profit and loss appropriation account?}
\sq{CLXIII}{What will be the entry for distribution of profit among the partners?}
\sq{CLXIV}{What will the entry for distribution of loss among the partners?}
\sq{CLXV}{What adjustments are necessary on the admission of a partner into partnership?}
\sq{CLXVI}{How a new partner can be entered into partnership?}
\sq{CLXVII}{What is sacrifice ratio?}
\sq{CLXVIII}{Define goodwill.}
\sq{CLXIX}{Write the three factors which are responsible for the goodwill of a business.}
\sq{CLXX}{Why the necessity of the valuation of goodwill arises?}
\sq{CLXXI}{How many methods are available for the valuation of goodwill, if the goodwill does not exist in the books of account?}
\sq{CLXXII}{What is average profits method for the valuation of goodwill?}
\sq{CLXXIII}{What is super profits method for the valuation of goodwill?}
\sq{CLXXIV}{Write the formula to calculate the value of goodwill according to capitalized method.}
\sq{CLXXV}{How many methods are available for the treatment of goodwill on the admission of a new partner, if goodwill does not exist in the books of account?}
\sq{CLXXVI}{What is premium method for the treatment of goodwill on the admission of a new partner?}
\sq{CLXXVII}{What is revaluation method for the treatment of goodwill on the admission of a new partner?}
\sq{CLXXVIII}{What is memorandum revaluation method for the treatment of goodwill on the admission of a new partner?}
\sq{CLXXIX}{Show the journal entries under premium method when the goodwill brought in cash and retained in the business.}
\sq{CLXXX}{Show the journal entries under premium method, when the goodwill brought in cash and withdrawn by the old partners.}
\sq{CLXXXI}{What journal entry will have to be made for goodwill raised under the revaluation method on the admission of a new partner?}
\sq{CLXXXII}{Pass the journal entries for goodwill under memorandum revaluation method of goodwill.}
\sq{CLXXXIII}{What will be the treatment of goodwill, if the goodwill account already exists in the books of the accounts?}
\sq{CLXXXIV}{What is revaluation account?}
\sq{CLXXXV}{Show the entries for increase and decrease in the value of assets on the admission of a partner.}
\sq{CLXXXVI}{What journal entries will have to be made for assets and decrease in the value of liabilities on the admission of a partner?}
\sq{CLXXXVII}{Pass the journal entries for revaluation profit or revaluation loss on the admission of a new partner.}
\sq{CLXXXVIII}{Distinguish between revaluation account and memorandum revaluation account.}
\sq{CLXXXIX}{Show the entry for distribution of reserve or distribution profit on the admission of a new partner.}
\sq{CXC}{Calculate the value of goodwill on 3 years purchase of the average profit of the last five years. The profits earned during the last five years are Rs. 6,500, Rs. 5,000, Rs 6,000, Rs.7,000 and Rs.10,000.}
\sq{CXCI}{Suppose the average profit is Rs.8000 and the reasonable return is 10\% the value of net assets is Rs.50,000.}
\sq{CXCII}{A \& B share profit and losses in the ratio of 3:2. They admit C into partnership. The new profit sharing ratio is 7:4:3. Calculate the sacrifice ratio of old partners?}
\sq{CXCIII}{A \& B share profit and losses in the ratio of 3:2. They admit C into partnership. The new profit sharing ratio is 7:5:4. Calculate the sacrifice ratio of old partners.}
\sq{CXCIV}{C \& D share profit and losses in the ratio of 3:1. They admit K into partnership. The new profit sharing ratio is 3:1:1. Calculate the sacrifice ratio of old partners.}
\sq{CXCV}{A \& B share profit and losses in the ratio of 3:2. They admit C into partnership. The new profit sharing ratio is 2:2:1. Calculate the sacrifice ratio of old partners?}
\sq{CXCVI}{A \& B share profit and losses in the ratio of 3:2. They admit C into partnership. The new profit sharing ratio is 5:3:2. Calculate the sacrifice ratio of old partners?}
\sq{CXCVII}{If C’s share is Rs.20,000 for 1/5 profit, then calculate the total capital in partnership?}
\sq{CXCVIII}{Write any three Accounting problems on retirement of a partner.}
\sq{CXCIX}{What will be the new profit sharing ratio, if a partner retires and there is no fresh agreement?}
\sq{CC}{What is the formula of gaining ratio or how gaining ratio can be calculated?}
\sq{CCI}{Distinguish between gaining ratio and sacrifice ratio. (any two)}
\sq{CCII}{Pass the journal entry, in case of retirement, when goodwill is raised at full value.}
\sq{CCIII}{Suppose the value of goodwill in the books at Rs.16000 and the value of goodwill on the retirement of a partner is estimated at Rs.36000, then pass the journal entry for goodwill.}
\sq{CCIV}{What journal entries will have to be made in case of retirement if the goodwill is raised at full value and written off?}
\sq{CCV}{What will be the journal entry, in case of retirement, if the goodwill is raised with retiring partner’s share only?}
\sq{CCVI}{Show the journal entries, in case of retirement, if goodwill is raised partially and written off.}
\sq{CCVII}{What journal entry will have to be made, in case of retirement, if the goodwill is not raised?}
\sq{CCVIII}{Suppose A \& B and C are partners sharing profits in the ratio of 3:2:1. Find out the new profit sharing ratio? If A retires, B retires, C retires, In the absence of any fresh agreement}
\sq{CCIX}{A, B and C are Partners in the ratio of 1/2, 3/10 and 2/10 respectively. B retires and his share is taken by A and C in the ratio of 2:1. Find out the new profit sharing ratio?}
\sq{CCX}{What is the treatment of Reserves and undistributed profit in case of retirement of a partner?}
\sq{CCXI}{What journal entry will have to be made, when the capital account of a retiring partner is transferred to loan account?}
\sq{CCXII}{Suppose A, B and C are equal partner, profit and loss account for the year ended December 31st 2015 revealed net profit Rs.36000. A dies on March 31st, 2016. Calculate the A’s share of profit based on previous year’s profit?}
\sq{CCXIII}{Suppose A, B and C are partners sharing profit in the ratio of 4:5:7. A dies on 31st March 2016. The sales and profit during 2015 were Rs.2, 40, 000 and 60,000 respectively. The sales upto March 31st, 2016 amounted to Rs.80, 000. Calculate A’s share in profit on sales basis?}
\sq{CCXIV}{Pass the journal entry for interest on capital of deceased’s partner.}
\sq{CCXV}{Pass the journal entry for accumulated profits / Reserves on the death of a partner.}
\sq{CCXVI}{What is the difference between dissolution of partnership and dissolution of the firm?}
\sq{CCXVII}{What is meant by dissolution of firm by agreement?}
\sq{CCXVIII}{What is compulsory dissolution or dissolution by the operation of law of a firm?}
\sq{CCXIX}{What is meant by dissolution by notice of firm at will?}
\sq{CCXX}{When a firm may be dissolved by court order? (any three points)}
\sq{CCXXI}{What is meant by dissolution of firm on the happening of certain contingencies?}
\sq{CCXXII}{What is realization account?}
\sq{CCXXIII}{Pass journal entry, for closing various assets account in case of dissolution of firm.}
\sq{CCXXIV}{Pass journal entry, for closing various liabilities account in case of dissolution of firm.}
\sq{CCXXV}{Pass journal entry, when assets are sold for cash in case of dissolution of firm.}
\sq{CCXXVI}{Pass journal entry, for payment of liabilities in case of dissolution of firm.}
\sq{CCXXVII}{Pass journal entry, for assets taken over by partner in case of dissolution of firm.}
\sq{CCXXVIII}{Pass journal entry, for payment of unrecorded liability in case of dissolution of firm.}
\sq{CCXXIX}{Pass journal entry, for sale of unrecorded asset in case of dissolution of firm.}
\sq{CCXXX}{Pass journal entry, for liabilities taken over by partner in case of dissolution of firm.}
\sq{CCXXXI}{Pass journal entry, for transfer of undivided profit in case of dissolution of firm.}
\sq{CCXXXII}{Pass journal entry, for expenses of realization in case of dissolution of firm.}
\sq{CCXXXIII}{Pass journal entry, for transferring loss on realization in case of dissolution of firm.}
\sq{CCXXXIV}{Pass the journal entry, for transferring profit on realization in case of dissolution of firm.}
\sq{CCXXXV}{Pass the journal entry, for repayment of partner’s capital balance in case of dissolution of firm.}
\sq{CCXXXVI}{Pass the journal entry, for repayment of partner’s loan in case of dissolution of firm.}
\sq{CCXXXVII}{What is Garner Vs Murray decision?}

\end{document}
